% SGA 3，第 VII_B 篇，英文重建稿——第 1 轮。
% 范围：权威底本分篇第 84--85 页，结束第 5.1.2 节的证明及推论，止于第 5.1.3 节之前。
% 合并阅读本页码： 597--598.
% 再版印刷页码： 586--587.
% 本范围内原版页边定位： 554.
% 本范围内编者注： 148.
% 权威底本： Polo--Gille Exp7B-13oct24.pdf,
% SHA-256 D303690E0195654BBE7C3DAA5D89EA23E24CD709A3575217234AC19AE8836694.
% 曾参考 Jacob Reinhold 的 CC BY 4.0 Markdown 作为比对谱系。
% 原英文 Loop 1 阶段说明：证明中各行正合的图采用直接取自 PDF 的精确裁片。当前中文工程已采用原生图。

\noindent\textit{证明。}
对每个 \(i,j\geq0\)，置
\[
  W_{i,j}=U_i\widehat\otimes V_j,
  \qquad
  \overline W_{i,j}
  =
  (U_i/U_{i+1})\widehat\otimes_k(V_j/V_{j+1}).
\]
我们对 \(n\) 归纳，证明自然映射
\[
  \pi_n:W_n\longrightarrow
  \bigoplus_{i+j=n}\overline W_{i,j}
\]
为满射，且包含关系
\[
  W_{n+1}\subset\ker(\pi_n)
\]
为等号。当 \(n=0\) 时，投影
\[
  \pi_0:
  U_0\widehat\otimes_kV_0
  \longrightarrow
  (U_0/U_1)\widehat\otimes_k(V_0/V_1)
\]
为满射。由于 \(U_0\)、\(U_0/U_1\)，从而 \(U_1\) 都在 \(k\) 上拓扑平坦，有
\[
  \ker(\pi_0)
  =
  U_0\widehat\otimes_kV_1
  +
  U_1\widehat\otimes_kV_0,
\]
并且
\[
  \bigl(U_0\widehat\otimes_kV_1\bigr)
  \cap
  \bigl(U_1\widehat\otimes_kV_0\bigr)
  =
  U_1\widehat\otimes_kV_1.
\]

现在假设 \(n>0\)，且已对 \(n-1\) 证明了结果。置
\[
  M_0=U_0\widehat\otimes_kV_n,
  \qquad
  S_0=\sum_{i=1}^{n}U_i\widehat\otimes_kV_{n-i}.
\]
有 \(S_0\subset U_1\widehat\otimes_kV_0\)。将上述论证用于 \(V_0\supset V_n\) 代替 \(V_0\supset V_1\)，得到
\[
  M_0\cap S_0
  \subset
  \bigl(U_0\widehat\otimes_kV_n\bigr)
  \cap
  \bigl(U_1\widehat\otimes_kV\bigr)
  =
  U_1\widehat\otimes_kV_n.
\]
由此可得
\[
  M_0\cap S_0=U_1\widehat\otimes_kV_n.
\]
由于 \(W_n=M_0+S_0\)，得到如下各行正合的交换图，其中
\[
  U'_i=U_{i+1},
  \qquad
  W'_{i,n-1-i}=W_{i+1,n-1-i}
  \quad(0\leq i\leq n-1):
\]
\begin{center}
  \phantomsection
  \label{sga3:VIIB:5.1.2:filtration-exact-rows-diagram}
  % SGA3-VII-DIAG-125
  % SGA3-VII-DIAG-125
% Authority: Exp7B-13oct24.pdf, local p. 85.
% Crop: figures/Exp7B_localp085_section512_filtration_exact_rows_diagram.png
% Locator: tex/components/78_expose_VIIB_section512_filtration_lemma_proof_corollary_en.tex, Loop 1 line 89.
\begin{tikzcd}[row sep=3em,column sep=1.8em]
0 \arrow[r] & S_0 \arrow[r] \arrow[d,"\pi'_{n-1}"']
  & W_n \arrow[r] \arrow[d,"\pi_n"']
  & (U_0/U_1)\widehat\otimes_k V_n
      \arrow[r] \arrow[d,"p"']
  & 0 \\
0 \arrow[r]
  & \displaystyle\bigoplus_{i=0}^{n-1}\overline W'_{i,n-1-i}
      \arrow[r]
  & \displaystyle\bigoplus_{i=0}^{n}\overline W_{i,n-i}
      \arrow[r]
  & \overline W_{0,n} \arrow[r]
  & 0
\end{tikzcd}

\end{center}

\phantomsection
\label{sga3:VIIB:5.1.2:diagram-source-notation-note}
\noindent\textit{底本说明。}\enspace
法文图中下行中间的项印作 \(\bigoplus_{i=0}^{n}\overline W_{i,n-1}\)；而紧接其上定义的自然映射 \(\pi_n\) 的目标为 \(\bigoplus_{i+j=n}\overline W_{i,j}\)，这将给出 \(\bigoplus_{i=0}^{n}\overline W_{i,n-i}\)。规范记录 \texttt{SGA-SRC-000840} 接受这个互补指标；上图采用此指标，同时在此保留印本的读法。

于是 \(p\) 为满射，其核为
\[
  (U_0/U_1)\widehat\otimes_kV_{n+1}.
\]
此外，将归纳假设用于序列 \((U'_i)\)，可知映射 \(\pi'_{n-1}\) 为满射，其核为
\[
  W'_n
  =
  \sum_{i=1}^{n}W_{i,n+1-i}.
\]
由此可知 \(\pi_n\) 为满射，且
\[
  W_{n+1}\subset\ker(\pi_n)
\]
为等号。引理得证。

回到 \(\mathcal C\) 的对象 \((A,J)\)。由 \SGARefLink{sga3:VIIB:0.2.G}{0.2.G}，每个
\[
  \overline{J^n}/\overline{J^{n+1}}
\]
都在 \(k\) 上拓扑平坦的假设，等价于 \(\operatorname{Gr}_J(A)\) 在 \(k\) 上拓扑平坦。

\medskip
\noindent\textbf{推论。}
\phantomsection
\label{sga3:VIIB:5.1.2:corollary}
{\itshape 令 \(\mathcal P\) 为 \(\mathcal C\) 的全子范畴，其对象为使 \(A\) 和 \(\operatorname{Gr}_J(A)\) 在 \(k\) 上拓扑平坦的对象 \((A,J)\)。则函子
\[
  \mathcal P\longrightarrow\mathbf{Alpg}/k,
  \qquad
  (A,J)\longmapsto\operatorname{Gr}_J(A),
\]
与有限余积交换，因而将 \(\mathcal P\) 的余群映至 \(\mathbf{Alpg}/k\) 的余群。

特别地，若 \(k\) 为域，则对 \(\mathcal C\) 的每个余群 \((A,J)\)，代数 \(\operatorname{Gr}_J(A)\) 为 \(\mathbf{Alpg}/k\) 的余群。}
\footnote{\label{sga3:VIIB:nde:148}因此，对形式 \(k\)-群对 \(H\subset G\)，可以联系一个“\(G\) 沿 \(H\) 的形式完备化 \(\widehat G_H\)”，它是形式 \(k\)-群。此外，我们将在 \SGARefLink{sga3:VIIB:5.1.3}{5.1.3}--\SGARefLink{sga3:VIIB:5.1.4}{5.1.4} 中看到，包含映射 \(\sigma:H\hookrightarrow\widehat G_H\) 有一个收缩映射 \(\pi:\widehat G_H\to H\)，而形式 \(k\)-群 \(N=\ker(\pi)\) 作为形式簇，与齐性空间 \(G/H\) 沿单位截面的完备化等同。这将在 \SGARefLink{sga3:VIIB:5.2.2}{5.2.2} 中用到。}
